\section{My support}

\textbf{4a. What others did that helped me work more effectively.} Four people changed how I worked, and I want to name the specific mechanism for each because a generic list of thanks is a Merit-ceiling move.

\textbf{Edward} is the reason the Cube ever spoke to the Pi. Around wk15 I had spent most of a weekend trying to make \texttt{pymavlink} read from \texttt{/dev/ttyAMA0} on Python 3.13, and it was Edward --- after I had filed three dead-end notes in Teams --- who pointed out that \texttt{pyserial} reads are broken on Python 3.13 and that the workaround was a \texttt{mavproxy} UDP bridge. That 40-minute conversation unlocked roughly two weeks of Pi testing that would otherwise have stalled. The uncomfortable realisation was not the technical fact; it was that my instinct to ``finish debugging before asking for help'' is a cost I pay, not a professionalism I exhibit. Edward's throwaway remark that ``just because the test passes on laptop doesn't mean the hardware does'' reframed how I thought about simulation fidelity for the rest of the project --- I stopped treating SITL parity as evidence and started treating it as a hypothesis.

\textbf{Robin} helped me by refusing to integrate against my first state machine. In wk17 I had a design I believed was technically correct and he pushed back until I could prove the extra states caused materially different decisions. He was right, I was partly wrong, and his pushback is the reason the \texttt{main.py} refactor (1411\,$\rightarrow$\,754 lines, 6 modules) happened at all. I now see that what I had framed as a technical disagreement was actually a maintainability disagreement, and I preferred the technical framing because it protected me from noticing I had built something only I could read.

\textbf{Demetro} helped me in a way that is easy to miss: he accepted drafts written in my voice without turning the collaboration into a status contest. Not every teammate would have done that, and the speed it unlocked on the FDR slides was entirely his generosity, not mine. \textbf{Our pilot} helped by operating the browser ground station unaided during the field-day dry runs despite wind and a cancelled flight --- the fact that he did not need to ask any questions validated the minimalist button layout and told me the UX was correct in a way no laptop testing could have. \textbf{Sid} (technical assistant, not teammate) delivered the single DJI flight video with SRT telemetry that I used for FOV calibration at nine altitudes, 300 synthetic backgrounds, 16 labelled real frames, and the A/B model comparison; the retrained model at mAP50\,=\,0.995 does not exist without that one video, and the honest lesson is that I should have asked for it three weeks earlier than I did.

\textbf{4b. What I did that helped others work more effectively.} I will put three specific feedback episodes on record, because the exemplar guide is explicit that ``I gave constructive feedback when needed'' is the single most common Merit-ceiling sentence in UK engineering reflection.

\textit{Feedback to Demetro --- FDR speaking scripts.} When Demetro asked for help with his two Final Design Review slides (commit \texttt{32fb0c5}) I built the layouts from scratch and then wrote full speaking scripts (246 + 219 words, 3.1\,min total, commits \texttt{b2a9cd4}, \texttt{ba667f2}) with an inline teleprompter sidebar, iterating 4--5 times (\texttt{9b826f8}, \texttt{7920e5a}, \texttt{794ef5c}). His first reaction was polite gratitude, but on the third pass he pushed back: the wording was not how he spoke. The observable behaviour change came on FDR day --- he delivered the two slides confidently, on time, in his own cadence, and a week later when he prepared his next set of slides himself he was noticeably using specific numbers rather than adjectives (a habit I had modelled in draft three). What I learned about my own delivery is harder than what he learned from mine: my first three drafts were in my register, not his, and I now see that good feedback to a peer presenter means writing in their voice, not polishing yours. The scarce resource was register, not prose quality.

\textit{Feedback to Robin --- CENTERING timeout and the handoff contract.} In wk19 I noticed Robin's flight script had no timeout guard on the \texttt{CENTERING} state: if the target was lost mid-descent the loop would hang on a detection that never came. I sent him a Teams message naming the exact variable (\texttt{consecutive\_lost}) and the exact line, and suggested a time-based timeout with fallback to ascend-and-retry. He fixed it the same afternoon --- the fix is now visible in \texttt{4\_detect\_and\_center.py}. The observable behaviour change came a week later: in wk20 he pulled me aside after the team meeting to say he had caught a similar timeout bug in a separate module because he was now ``looking for that shape of bug.'' That is the outcome I wanted --- not the specific patch, but the shape of the check propagating into his working model. Around the same time I wrote \texttt{passive\_watch\_2.py} with a \texttt{--robin} flag, an HTTP \texttt{cv\_mode} polling endpoint, and a dedicated README with the JSON schema in copy-pasteable form (commits \texttt{a5b1ab7}, \texttt{885cee4}, \texttt{36b6069}, \texttt{24783d9}). Robin integrated the vision stack in a single session without asking follow-up questions for two days. The behaviour change that matters most to me is that he started writing similar caller-shaped contracts for \emph{his} handoffs to other teammates --- the habit propagated. The generalisable lesson is that when giving feedback to a high-performing peer, specificity is the scarce resource, not softness.

\textit{Feedback to the pilot --- ground-station button layout.} When I showed the pilot the first \texttt{pi\_flight.py} dashboard he said the button grid was too dense to hit reliably while holding an RC transmitter. I proposed a minimal four-button set (Y / N / M / L) mapped to the RC thumb-reach pattern, and wired an emergency RTL into the override guard so mode 9 (LAND) from the RC switch would be honoured. During the field-day dry runs he triggered the abort three times without looking at the screen. The observable behaviour change I did not expect: on the same day he asked for a similarly minimal layout on the VERIFY confirmation screen, which I had designed more densely. That request is the point at which my frame shifted --- I had been designing for me-as-operator in a chair, and the feedback loop forced me to see that the best UI is the one the pilot does not have to notice.

One concrete register-translation episode is worth naming: when I filed the ``ACK race in \texttt{MAV\_CMD\_NAV\_TAKEOFF}'' bug report, our PM could not use it in her weekly Flight-Lab status write-up, so I rewrote it as ``the drone doesn't lift off because it is waiting for a confirmation message that never comes,'' and she escalated \emph{that} plain-language version to Steve --- who actioned it the same day. The lesson is not that simpler words are better; it is that the register I default to protects me from audiences I need to reach.

\textbf{A/B methods comparison --- handoff docs vs in-person walkthrough.} I tried two approaches to the same problem of onboarding Robin onto the vision stack: (A) a 3000-word Markdown colleague handoff document written alone over an evening, and (B) a 45-minute in-person pair walkthrough on the same material. Measured against time-to-first-successful-integration, (A) scored $\sim$10 days (he finally read it when forced to), (B) scored under an hour. (B) was more effective because it forced me to notice which parts of my mental model were load-bearing for a reader and which were decoration, and it let him interrupt. The forward-looking rule: any document longer than 500 words must be paired with a live walkthrough inside a week, or it should be treated as a draft for \emph{me}, not an artefact for the team.

The non-feedback infrastructure I built for the team's use was the quieter half of this work: \texttt{brain-dump.md}, \texttt{MEMORY.md}, \texttt{NICE\_TO\_HAVE.md} (institutional memory that survived session boundaries), the blueprint system for files longer than 500 lines (\texttt{docs/main\_blueprint.md} and three others), \texttt{CLAUDE.md} as project ground truth for new sessions, the \texttt{simple\_simulator.py} web ground station that let teammates test without the drone, and \texttt{FIELD\_QUICK\_REF.md} for no-internet field use. I also drafted the formal clarification letter to the course organiser (see Appendix, \texttt{docs/COMPLAINT\_LETTER.md}) --- a non-technical-audience artefact the whole team reviewed and signed, and the piece of writing that taught me most about register, because my first three drafts were in my voice and my anger, and the fourth was in the team's voice only after Robin cut three sentences that were ``technically accurate but personally sharp.'' He was right.

The uncomfortable closing realisation is this. I now see that building tools for others was my version of teamwork \emph{because} I struggled with real-time collaboration --- async handoff docs were emotional scaffolding for me as much as they were help for the recipient. My next project must invest in the harder skill of pairing, even though pairing feels less productive per hour, because the peer-development outcomes that matter most (Demetro's cadence, Robin's bug-shape instinct, the pilot's minimalist instinct) happened in the few moments I \emph{did} collaborate synchronously, not in any of the documents I wrote alone.

\textbf{Closing: identity shift.} I came into this project believing that engineering excellence was a solo output: the thesis that ``simulation-first delivery plus 820 commits equals a good project'' was really a thesis about who I thought I had to be. I am leaving it with a different identity, and I want to state it in falsifiable terms. I no longer believe that the engineer who ships alone is the engineer who matters most; I believe that the engineer whose teammates become more capable \emph{because of the interaction} is the one who matters, and the evidence I trust is behavioural --- Robin's bug-shape instinct, Demetro's numbers-not-adjectives cadence, the pilot's minimalist UI instinct. That reframes Ukrainian engineering culture for me too: the instinct to ``just fix it yourself'' is survival pragmatism from a scarcity context, not a virtue to export into team projects in slack-rich environments. The falsifiable signal for the next twelve months: on my next project, I will log every synchronous pairing session and every solo sprint, and the ratio will move. If at month twelve pairing hours are still under 20\% of build hours, I will have failed to change, and I will know it by the number, not by how I feel about it.
